\chapter{知识的融合}
\label{chap:knowledge-fusion}

同一张花卉照片可能得到几份不同答案：人工认为是茶梅，两条同模板规则都报山茶，模型检查后没有作答。我们已经取得了多条知识记录，却还不知道该采用哪个类别。两张相同的规则票可能重复了一份证据，人工也可能出错；若记录实际标注的是照片中的不同花朵，连“答案冲突”都尚未成立。本章从这些记录出发，研究怎样判断哪些证据可以合用，怎样形成综合结果，以及结果还应保留哪些不确定性。

要解释这次分歧，我们需要知道来源在哪些样本上作答、会怎样出错，再确认记录是否指向同一对象和同一命题，最后比较组合证据的规则。正文按来源建模、知识对齐、融合推断和结果使用展开；这是讲解顺序，实际利用标签估计来源参数之前必须完成相应对齐，有些参数还要与未知标签共同估计。本章沿用第~\ref{sec:knowledge-acquisition-problem}节的知识记录与状态约定，并使用前卷的条件概率、极大似然和损失函数；新的来源假设会在使用处说明。证据不足时可以保留概率、候选集合或“未决”状态，由使用端决定是否弃权；不能回答同一命题的关系与恢复目标，则保留各自含义联合使用。

\section{知识来源的建模}
\label{sec:knowledge-fusion-source-modeling}

面对人工、规则和模型给出的标签，我们先把每份判断追溯到产生它的过程。\term{知识来源}（knowledge source）就是这个产生知识记录的主体或过程，可以是一名标注者、一条规则、一个模型，或模型的一种固定配置。两列输出若共用模板、词典、训练数据或基础模型，仅有不同列名不足以说明它们独立。来源模型因而同时读取来源身份、覆盖状态、已对齐输出和共享关系，用这些记录描述能力与依赖，为后续证据组合提供参数。

同一份来源记录需要从三个方面解释。\term{可靠性}（reliability）在本章的标签任务中描述来源作答后有多大比例答对；\term{混淆}（confusion）描述真实类别给定时，它会报成哪一类；\term{相关性}（correlation）考察多个来源的输出或错误是否共同变化，其中共享信息造成的依赖尤其会影响融合。正确比例不能说明错误方向，分别估计每个来源也不能说明联合行为。两个总体正确率相同的来源可能错向不同类别，四个高正确率来源也可能只是同一错误机制的四份副本。

\subsection{来源可靠性的建模}
\label{sec:knowledge-fusion-reliability}

若一个来源只回答容易照片，把未作答与答错合成一个指标就无法解释它的能力。我们先分开“是否有答案”和“答案是否正确”：对$n>0$个已完成对象与语义对齐的样本，设有$m$个来源，共同类别集合为$Y=\{1,\ldots,K\}$；来源$s$在样本$i$上的有效类别记为$\lambda_{i,s}\in Y$，有有效类别时$o_{i,s}=1$，否则为$0$。这个掩码只为计算选出可比较标签；计算时先筛出有效位置，再读取类别、概率或监督内容，不给无效位置填零类别或零目标。原记录仍保留执行状态、来源动作和取得条件。例如规则运行成功但未触发时，动作为空、内容未知，触发结果另记，不能把它改记为一次主动弃权。

\subsubsection{利用参考标签估计可靠性}
\label{sec:knowledge-fusion-reference-reliability}

如果部分照片已有独立审计形成的参考标签，我们可以核对来源的答案，而不必先从来源赞同关系猜测真值。记审计索引集合为$I_{\mathrm{audit}}$、参考标签为$y_i$。当前需要两个可计算量：覆盖率衡量来源在全部目标样本中回答了多少，覆盖条件下正确率衡量它在有参考且实际回答的样本中答对多少。两者分别按各自样本集合取比例：
\begin{equation}
  \widehat c_s
  =\frac{1}{n}\sum_{i=1}^{n}o_{i,s},
  \qquad
  \widehat a_s
  =\frac{\sum_{i\in I_{\mathrm{audit}}:o_{i,s}=1}
    \mathbb{I}\{\lambda_{i,s}=y_i\}}
  {\sum_{i\in I_{\mathrm{audit}}}o_{i,s}}.
  \label{eq:knowledge-fusion-audited-reliability}
\end{equation}
覆盖率的分母是$n$，没有答案的目标样本仍占一个位置；正确率的分母则只数审计集合中$o_{i,s}=1$的样本，分子再数其中与参考一致的答案。第二个分母为零时，正确率尚未估计，不能填成零或一。未请求、失败和弃权会降低覆盖，却不直接构成类别错误；若把弃权计作负类，统计的已不再是原来的分类任务。

表~\ref{tab:knowledge-fusion-source-records}把这两个分母放在同一组教学记录中。规则$R_1,R_2$共享模板，人工来源$H$独立作答，模型$M$可以弃权；前四张照片有独立参考，后四张没有参考。我们先用前四行核对正确性，再用全部八行计算覆盖，后文仍用这些来源输出比较融合规则。表中记录均为人工构造，不是实验结果。

\begin{table*}[htbp]
  \centering
  \small
  \renewcommand{\arraystretch}{1.15}
  \begin{tabular}{@{}clccccc@{}}
    \toprule
    \textbf{样本} & \textbf{区域} & \textbf{参考} & \textbf{$R_1$} & \textbf{$R_2$} & \textbf{$H$} & \textbf{$M$} \\
    \midrule
    $x_1$ & 简单 & 山茶 & 山茶 & 山茶 & 山茶 & 山茶 \\
    $x_2$ & 困难 & 茶梅 & 山茶 & 山茶 & $\varnothing_{\text{未请求}}$ & 茶梅 \\
    $x_3$ & 简单 & 山茶 & 山茶 & 山茶 & 山茶 & 茶梅 \\
    $x_4$ & 困难 & 茶梅 & $\varnothing_{\text{失败}}$ & $\varnothing_{\text{失败}}$ & $\varnothing_{\text{未请求}}$ & 山茶 \\
    \midrule
    $x_5$ & 简单 & 未审计 & 山茶 & 山茶 & 茶梅 & $\bot$ \\
    $x_6$ & 简单 & 未审计 & 山茶 & 山茶 & 茶梅 & 茶梅 \\
    $x_7$ & 未覆盖 & 未审计 & $\varnothing_{\text{失败}}$ & $\varnothing_{\text{失败}}$ & $\varnothing_{\text{未请求}}$ & $\bot$ \\
    $x_8$ & 简单 & 未审计 & 茶梅 & 茶梅 & 山茶 & 山茶 \\
    \bottomrule
  \end{tabular}
  \caption{多来源标签与覆盖记录。$\varnothing_{\text{未请求}}$表示没有发起请求，$\varnothing_{\text{失败}}$表示执行失败，不能用来代指仍在等待或处理的任务；$\bot$表示成功执行后主动弃权，内容仍为未知。这三种情况都不计类别票。$x_1$至$x_4$有独立参考，用于教学与开发诊断，不充当最终留出评价；所有对象和类别语义均假设已经对齐。}
  \label{tab:knowledge-fusion-source-records}
\end{table*}
\FloatBarrier

逐行核对表~\ref{tab:knowledge-fusion-source-records}，$R_1,R_2,M$的有效覆盖率都是$6/8$，仅在简单区作答的$H$为$5/8$。审计区中，两条规则各答对$2/3$，$H$答对$2/2$，$M$答对$2/4$。这些比例的分母不是同一批对象：$H$没有回答困难区。把比较限制到共同覆盖的简单审计样本$x_1,x_3$，$H$为$2/2$、$M$为$1/2$；困难区$x_2,x_4$只有$M$的$1/2$，不能估计$H$。所以本表支持的是$H$在这两个简单样本上答得更准，而非$H$在全范围更可靠。区域$g$的正确率要把式~\eqref{eq:knowledge-fusion-audited-reliability}的分子、分母同时限制到$I_{\mathrm{audit}}\cap I_g$，并报告有效分母；更细分组也意味着更多审计投入。

这样的核对仍以参考标签和抽样范围为条件。参考答案也可能有误，审计的抽取方式、形成过程和规范版本都要对应目标使用范围；两次答对不能当作精确的能力估计。参考少时应保留区间或参数后验，无参考区域记为未估计。对选择性作答的来源，还要检查$o_{i,s}$随类别、难度和时间怎样变化，才能分清正确率升高来自判断改善，还是来自只留下容易题。

\subsubsection{利用来源关系估计可靠性}
\label{sec:knowledge-fusion-relational-reliability}

现在去掉参考标签，只留下三列始终作答的二分类记录。假设一半样本三列都报山茶，另一半都报茶梅，我们能确认来源一致，却不能直接核对正确性。解释甲把三列视为全部正确，两类各占一半；解释乙把三列视为全部报反，两类也各占一半。两种解释产生相同的观测分布，对每个样本的真值却作出相反判断。

增加同类记录也不能区分这两种解释，因为差异不在观测数量，而在观测根本没有确定类别方向。\term{可识别性}（identifiability）就是检查不同参数解释能否由观测分布区分；若两组参数始终给出同一分布，数据不能唯一确定参数。少量独立参考可以给类别命名，“来源倾向于报对”的假设也可以排除报反解释，但这时排除依据是新增约束，不是答案已经得到验证。

从记录估计能力时，可见量是标签矩阵$\bm{\Lambda}$和覆盖掩码$\bm{O}$，未知量是来源参数、类别比例和潜在答案。共同样本上的赞同与分歧把这些未知量联系起来，但是否能恢复参数还取决于重叠和信息差异。两个来源从不共同覆盖，且没有其他重叠来源连接它们时，缺少比较的桥梁；十列逐项复制同一规则时，又没有十份独立证据。缺少重叠、复制信息和类别置换是不同障碍，不能都用增加样本数解决。

要继续估计，我们可以减少待估参数、合并重复来源或设置明确先验，但每种处理都改变了允许的解释，不能代替真值核验。第~\ref{sec:knowledge-fusion-latent-label-inference}节将用来源生成标签的模型，把后验计算与参数更新接起来，并分析排除额外歧义的充分条件。更新能否停止是数值问题，观测能否区分参数则是可识别性问题。

\subsection{来源混淆与偏差的建模}
\label{sec:knowledge-fusion-confusion}

在三分类花卉任务中，两个来源都答对八成照片，仍可能对同一次分歧提供不同证据：一个常把茶梅报成山茶，另一个常把茶梅报成月季。总体正确率没有保留这种方向。为判断“若真值是$k$，当前输出$\ell$有多可能”，我们需要来源$s$的混淆矩阵；以$Y_i$表示样本$i$的真实类别，矩阵逐行记录给定$Y_i=k$且来源作答时，各个输出类别的条件概率：
\begin{equation}
  \theta^{(s)}_{k,\ell}
  =\mathbb{P}(\lambda_{i,s}=\ell
    \mid Y_i=k,o_{i,s}=1),
  \qquad
  \sum_{\ell=1}^{K}\theta^{(s)}_{k,\ell}=1.
  \label{eq:knowledge-fusion-confusion-matrix}
\end{equation}
行索引$k$固定真实类别，列索引$\ell$遍历共同集合$Y$中的输出，因而一行概率之和为一。对角项给出该类的正确率，非对角项给出错误流向。后文看到输出$\ell$时，将比较各候选真值对应行的同一列；实现和图表必须沿用这个行列方向。

已有参考时，我们在审计集合内只数来源实际作答的记录。参考为$k$、输出为$\ell$的次数记为$N^{(s)}_{k,\ell}$，在非负且行和为一的概率中最大化这些记录的似然，得到$\widehat\theta^{(s)}_{k,\ell}=N^{(s)}_{k,\ell}/\sum_rN^{(s)}_{k,r}$。分母是该真实类别下的有效作答数，零分母表示这一行未估计。少数类别计数很少时，可以使用Dirichlet平滑，即为整行概率加入明确的先验后再估计，或合并过细类别、报告区间；平滑缓解零计数，却不能新增真实观测。

回到这两个来源，设它们在同一目标范围上的总体正确率均为$0.8$，A主要把茶梅报成山茶，B主要把茶梅报成月季。当A报山茶而B报月季时，茶梅仍可能解释这组分歧：需要读取A矩阵的“茶梅、山茶”项与B矩阵的“茶梅、月季”项，再在相应依赖假设下组合。两个$0.8$无法完成这个计算，也不能单独确定最终类别。这里的比例和错误方向都是教学设定，只用于解释为何需要分类别似然。

输出“山茶”的比例升高，还不能定位是哪一部分发生变化。来源更多地选择山茶样本作答，改变的是覆盖机制；在相同真值下更偏向山茶，改变的是混淆矩阵；目标数据中山茶增多，改变的则是类别先验$\mathbb{P}(Y=k)$。这三种变化都能影响输出比例，所以仅观察边缘分布不足以认定来源退化，需要结合覆盖与分类别参考。

一张全局矩阵把所建模范围内的来源行为视为稳定。如果困难照片与简单照片上的混淆不同，就需要分组参数$\theta^{(s,g)}_{k,\ell}$，或让概率依赖样本特征；原来的行列含义不变，估计范围变小了。分组越细，参数越多，每行获得的参考与重叠也越少，应按使用需要划组并核对有效样本量和更新频率。若要让弃权本身参与推断，可以把它作为额外输出状态，或单独建模覆盖；两种方式都不把弃权当作真实物种类别。

\subsection{来源相关性的建模}
\label{sec:knowledge-fusion-correlation}

表~\ref{tab:knowledge-fusion-source-records}中的$R_1,R_2$常同时报山茶，但它们共用模板，第二张票是否新增信息需要另查。共享训练集、基础模型、检索结果、提示模板或标注指南也会造成类似问题。我们先把这些已知关系写入来源目录，再用已对齐输出检查是否还有未描述的依赖，而不从列数直接推断证据数。

一致率高可以来自两种机制：两个来源都随真实类别改变，也可以在真实类别已固定后仍一起出错。只有固定真值后，才能检查分别描述每个来源是否足够。对始终作答的$s,t$，若联合输出概率等于两个条件概率的乘积，就称给定$Y_i$后条件独立：
\begin{equation}
  \mathbb{P}(\lambda_{i,s},\lambda_{i,t}\mid Y_i)
  =
  \mathbb{P}(\lambda_{i,s}\mid Y_i)
  \mathbb{P}(\lambda_{i,t}\mid Y_i).
  \label{eq:knowledge-fusion-conditional-independence}
\end{equation}
这条等式正是把两份来源似然相乘的依据，不要求未给定真值时两列输出独立。等式不成立时，仅有各自混淆矩阵无法还原联合误差。没有参考真值而借助潜在标签检查的“残余相关性”，即扣除当前模型所解释的类别影响后剩余的共同变化，仍受模型假设限制；一个经验相关系数不能唯一确定共享机制或因果方向。

图~\ref{fig:knowledge-fusion-source-dependence}把这一区别画成两种假设。左侧每份输出只按各自条件分布响应真实类别，右侧三条规则还受同一模板或证据影响。我们应观察右侧的虚线和来源组：它们提示三张赞同票可能应作为一组证据处理，而不是自动计作三份独立判断。

\begin{figure*}[htbp]
  \centering
  \begin{tikzpicture}[
  x=1mm,y=1mm,
  every node/.style={font=\figurefont\small,text=BookInk},
  truth/.style={circle,draw=BookTeal,fill=FigureModel!12,line width=1pt,minimum size=11mm,inner sep=1pt},
  source/.style={book node,draw=BookBlue,fill=FigureInput!9,line width=0.8pt,minimum width=18mm,minimum height=9mm,inner sep=2pt},
  shared/.style={book node,draw=BookAmber,fill=BookAmber!10,line width=0.8pt,minimum width=24mm,minimum height=9mm,inner sep=2pt}
]
  \node[font=\figurefont\bfseries] at (35,26) {给定类别后独立};
  \node[truth] (yl) at (35,10) {$Y$};
  \node[source] (l1) at (8,-10) {来源1};
  \node[source] (l2) at (35,-10) {来源2};
  \node[source] (l3) at (62,-10) {来源3};
  \draw[book arrow,knowledge arrow,draw=BookTeal,line width=1pt] (yl) -- (l1);
  \draw[book arrow,knowledge arrow,draw=BookTeal,line width=1pt] (yl) -- (l2);
  \draw[book arrow,knowledge arrow,draw=BookTeal,line width=1pt] (yl) -- (l3);
  \node[font=\figurefont\scriptsize,text=BookMuted,align=center] at (35,-26) {三条标注路径分别提供类别票};

  \draw[BookRule,line width=0.7pt] (73,-40) -- (73,29);

  \node[font=\figurefont\bfseries] at (124,26) {共享机制引入相关错误};
  \node[truth] (yr) at (124,10) {$Y$};
  \node[source] (r1) at (86,-10) {规则1};
  \node[source] (r2) at (109,-10) {规则2};
  \node[source] (r3) at (132,-10) {规则3};
  \node[source] (human) at (155,-10) {标注者$H$};
  \node[shared] (u) at (109,-38) {共享模板或证据};
  \draw[book arrow,knowledge arrow,draw=BookTeal,line width=1pt] (yr) -- (r1);
  \draw[book arrow,knowledge arrow,draw=BookTeal,line width=1pt] (yr) -- (r2);
  \draw[book arrow,knowledge arrow,draw=BookTeal,line width=1pt] (yr) -- (r3);
  \draw[book arrow,knowledge arrow,draw=BookTeal,line width=1pt] (yr) -- (human);
  \draw[book arrow,knowledge arrow,draw=BookAmber,dashed] (u.north west) -- (r1.south);
  \draw[book arrow,knowledge arrow,draw=BookAmber,dashed] (u.north) -- (r2.south);
  \draw[book arrow,knowledge arrow,draw=BookAmber,dashed] (u.north east) -- (r3.south);
  \draw[BookAmber,dashed,line width=0.8pt,rounded corners=1.2mm] (75,-21) rectangle (143,0);
  \node[font=\figurefont\scriptsize,text=BookAmber,anchor=south west] at (78,3) {同一来源组};
  \node[font=\figurefont\scriptsize,text=BookAmber,anchor=west] at (78,-28) {共享影响};
\end{tikzpicture}

  \caption{独立信息与共享错误来源。左面板表示给定真实类别后来源独立的模型；右面板包含三条同模板规则与一名独立标注者，虚线和分组框标出三条规则的共享影响。图展示两种待检验的模型假设，不表示可以只从输出一致率唯一恢复真实依赖结构。}
  \label{fig:knowledge-fusion-source-dependence}
\end{figure*}
\FloatBarrier

为比较不同处理，我们固定右面板的输出：三条规则都报山茶，$H$报茶梅。逐列计票给山茶$3:1$的多数；若三条规则完全重复，合成一组后只剩山茶、茶梅各一票，结果未决。若还知道规则组与$H$的能力，就可以比较两类产生这组观察的概率，得到不同于票数的支持程度。

假设两类先验各为$0.5$，规则组与$H$始终作答且给定真值后独立，两者对称二分类正确率分别为$0.7,0.8$。若真值是茶梅，规则组报错而$H$报对，连同先验得到权重$0.5\times0.3\times0.8=0.12$；若真值是山茶，则得到$0.5\times0.7\times0.2=0.07$。把茶梅权重除以两类权重之和，得到后验$\mathbb{P}(Y=\text{茶梅}\mid\bm{\lambda})\approx0.632$，即当前模型看到这些输出后分给茶梅的概率。总体正确率能代入，是因为假设两类错误对称；错误不对称时必须读取混淆矩阵。这里三种结果固定了同一观察，只改变重复信息和能力的处理，尚未用独立真值证明哪种更准确。

完全复制时合并来源可以避免重复计票；有共享信息也有独有信息时，可以先给来源组分配总权重，再分配组内贡献；证据充足时，还可以在联合模型中加入成对或高阶依赖因子。这里的因子是描述多个输出共同作用的函数，后文将给出形式。合并过度会丢掉差异，细化依赖则增加参数与推断成本。弱监督的依赖结构学习也受稀疏性、样本量和模型条件限制，不能把可学习某类依赖扩成无标签数据能恢复任意关系的保证\scite{varma2019dependencies}

\section{多源知识的对齐}
\label{sec:knowledge-fusion-alignment}

两份记录分别写着“人物”和“组织”，我们还不能据此判定来源意见冲突：它们可能标注了文本中的不同实体，也可能采用不同任务规范。来源建模描述谁怎样作答，知识对齐则利用对象定位、内容、来源、规范和时间版本，确定哪些记录能够一起解释。对齐输出包括对应映射及其状态，可以是确定对应、部分对应、多个候选、待复核或无法转换。投票和概率标签模型只直接接收已确定落入共同答案空间的记录，不能为凑齐列数强造对应。

\subsection{知识对象的对齐}
\label{sec:knowledge-fusion-object-alignment}

先看同一句话中的“林舟”和“海岬科技”：人物标签若落在前者、组织标签落在后者，两者可以同时正确。\term{对象对齐}（object alignment）要确定不同记录究竟指向哪一个样本、实体、跨度、区域或时间区间。若对象标识在同一命名范围和版本内唯一且一致，可以据此精确对应；没有共享标识时，还要比较边界、位置、时间与上下文。关系记录需要同时对齐两端和方向，“海岬科技任命林舟”与“林舟任命海岬科技”不能因实体相同就合并。

我们可以先按精确标识匹配，再核对边界或时间重叠，仍有多个候选时才比较内容相似度；每一步都保留原始坐标、依据和确定程度。整段标注可能对应多条局部记录，时间区间也可能仅部分重叠，这时输出一对多或部分对应比强制选一条更准确地表达现有证据。扩大搜索范围让更多真实对应进入候选，也让更多同名对象参与竞争，同时增加匹配计算量。

图~\ref{fig:knowledge-fusion-object-alignment}固定同一句短文本，只改变两个标签的对象对应。左侧虚线把不同实体配在一起，制造了虚假冲突；右侧两条记录都指向“林舟”，此时“人物”和“组织”才是同一对象上的类别分歧。关系箭头另存端点与方向，承担的不是实体类别判断，不能混入这两张类别票。

\begin{figure*}[htbp]
  \centering
  \begin{tikzpicture}[
  x=1mm,y=1mm,
  every node/.style={font=\figurefont\small,text=BookInk},
  token/.style={book node,draw=BookRule,fill=white,minimum height=8mm,inner sep=2pt},
  span a/.style={draw=BookBlue,fill=FigureInput!10,line width=0.9pt,minimum height=7mm,inner sep=1.5pt},
  span b/.style={draw=BookAmber,fill=BookAmber!10,line width=0.9pt,minimum height=7mm,inner sep=1.5pt}
]
  \node[font=\figurefont\bfseries] at (36,28) {对象错配造成伪冲突};
  \node[token,minimum width=14mm] (ll1) at (9,13) {林舟};
  \node[token,minimum width=9mm] (join1) at (24,13) {在};
  \node[token,minimum width=13mm] (sh1) at (37,13) {上海};
  \node[token,minimum width=9mm] (verb1) at (49,13) {加入};
  \node[token,minimum width=22mm] (co1) at (66,13) {海岬科技};
  \node[span a,minimum width=20mm] (a1) at (9,-2) {来源A：人物};
  \node[span b,minimum width=26mm] (b1) at (66,-16) {来源B：组织};
  \draw[BookBlue,line width=0.9pt] (ll1.south) -- (a1.north);
  \draw[BookAmber,line width=0.9pt] (co1.south) |- (b1.north);
  \draw[book arrow,knowledge arrow,draw=BookAmber,dashed] (a1.south east) -- node[below,font=\figurefont\scriptsize,text=BookAmber] {错误对应} (b1.north west);
  \node[font=\figurefont\scriptsize,text=BookMuted,align=center] at (36,-29) {标签落在不同实体上\\不能作为同一对象的类别冲突};

  \draw[BookRule,line width=0.7pt] (80,-35) -- (80,32);

  \node[font=\figurefont\bfseries] at (121,28) {对齐后识别真实分歧};
  \node[token,minimum width=14mm] (ll2) at (94,13) {林舟};
  \node[token,minimum width=9mm] (join2) at (109,13) {在};
  \node[token,minimum width=13mm] (sh2) at (122,13) {上海};
  \node[token,minimum width=9mm] (verb2) at (134,13) {加入};
  \node[token,minimum width=22mm] (co2) at (151,13) {海岬科技};
  \node[span a,minimum width=20mm] (a2) at (94,-2) {来源A：人物};
  \node[span b,minimum width=20mm] (b2) at (94,-16) {来源B：组织};
  \draw[BookBlue,line width=0.9pt] (ll2.south) -- (a2.north);
  \draw[BookAmber,line width=0.9pt] (a2.south) -- (b2.north);
  \draw[BookTeal,line width=1pt,knowledge arrow] (ll2.south east) to[out=-25,in=-155] node[below,font=\figurefont\scriptsize,text=BookTeal] {加入：主体$\to$组织} (co2.south west);
  \node[font=\figurefont\scriptsize,text=BookMuted,align=center] at (128,-29) {同一跨度上的“人物/组织”\\才是需要融合的类别分歧};
\end{tikzpicture}

  \caption{对象错配与真实标签冲突。上、下两行分别表示两个来源的标注跨度；左面板的虚线把不同实体错误配成同一对象，右面板先用实线对齐“林舟”，再显示“人物”与“组织”的类别分歧。关系箭头单独标明两端和方向。}
  \label{fig:knowledge-fusion-object-alignment}
\end{figure*}
\FloatBarrier

标识之外的匹配需要保留疑义：同名对象可能不同，同一实体在不同时刻的状态也可能不同。若以相似度阈值接收对应，应同时核对漏配与错配，并复核阈值附近的样本。映射正确只说明找到了同一对象；两条准确对应的记录仍可能一起标错，因此对象对齐成功率不能替代标签正确率。

\subsection{知识语义的对齐}
\label{sec:knowledge-fusion-semantic-alignment}

对象已对应以后，同名标签仍可能回答不同问题。一条“正类”记录可以表示文本提到某项内容，另一条却表示该内容在现实中成立。\term{语义对齐}（semantic alignment）就是核对来源的判定规范、任务条件和有效版本，确定是否为同一命题，并说明类别怎样对应。标签名称是查找线索，不是含义相同的充分依据；对不同命题增加票数，不能解决定义差异。

为保存这种对应，设来源类别集合为$Y_s$、目标共同空间为$Y^\star$，用$h_s:Y_s\to 2^{Y^\star}$把来源类别映到目标候选集合；$2^{Y^\star}$表示$Y^\star$的全部子集。定义相同的类别可以映到单元素集合，上位类别只能限定若干细类候选；定义交叠而互不包含时，还要保留交叠条件与无法映射的部分。记录是否可逆、是否依赖上下文和规范版本，才能知道转换结果能否还原，以及映射修订影响了哪些旧记录。

把映射用于记录时，先转换类别表或候选集合，再决定是否进入融合。例如A的“相关”只要求出现关键词，B的“相关”要求文本证实命题，A的输出最多为B提供证据，不能直接算作B的一票。无预设语义的簇编号和表示向量也有相同障碍：编号相同或维度相同不建立语义对应，必须借助生成机制、已知对应的锚点样本或下游目标。

规范修订还会改变可比较范围。旧类别拆成两个新细类时，旧记录通常只约束两类组成的集合，不能反推出具体细类；原来的对应关系不再足以支持新任务。保存原始类别、规范版本和转换链，并在迁移时重新核对映射，才能区分旧知识仍支持的集合判断与尚需新增的细类证据。若只统一编码，概率模型再精细也会组合错误命题。

\subsection{知识粒度与形式的对齐}
\label{sec:knowledge-fusion-granularity-form-alignment}

两条记录即使含义相容，也可能一个只给父类，另一个给细类概率。\term{粒度对齐}（granularity alignment）处理类别层级或对象范围粗细不同的情形；\term{形式对齐}（form alignment）处理硬标签、概率、分数、集合等不同表达。我们要确定允许哪些转换，以及转换后还知道什么。子类互斥且穷尽父类时，子类概率可以相加得到父类概率；反向只有父类标签时，却没有依据把质量分给某个细类。

对象范围的转换也要核对任务语义。若文档正类定义为“至少出现一段正向证据”，局部正例可以推出文档正类；换成“全文支持该命题”，同样的局部记录就不够。概率转硬标签需要接收阈值或决策损失，任意分数转概率需要尺度与校准依据；这里的校准是检查输出概率能否对应真实事件的频率。没有调用、调用失败和运行后弃权虽然都可能没有类别值，其生成条件不同，转换时仍须分开保存。

表~\ref{tab:knowledge-fusion-form-conversions}逐行列出转换条件、结果保留的信息和无法恢复的信息。不可逆转换仍可满足接口需要，但应留下原记录与配置，以便追踪在哪一步丢掉了细类或不确定性。若根本没有进入共同答案空间的合理映射，就保留原监督对象，交由第~\ref{sec:knowledge-fusion-heterogeneous-use}节的联合训练安排使用。

\begin{table*}[htbp]
  \centering
  \small
  \renewcommand{\arraystretch}{1.15}
  \begin{tabularx}{\linewidth}{@{}>{\raggedright\setlength{\parindent}{0pt}\arraybackslash}p{28mm}>{\raggedright\setlength{\parindent}{0pt}\arraybackslash}X>{\raggedright\setlength{\parindent}{0pt}\arraybackslash}X>{\raggedright\setlength{\parindent}{0pt}\arraybackslash}X@{}}
    \toprule
    \textbf{转换} & \textbf{操作与条件} & \textbf{保留信息} & \textbf{丢失或待补信息} \\
    \midrule
    细类到粗类 & 合并互斥且穷尽的子类，概率求和 & 父类归属及总概率 & 原细类及其概率分配 \\
    粗类到候选细类 & 保留父类下全部合法子类 & 真值所属父类 & 具体子类，需新增证据 \\
    概率到硬标签 & 按阈值或决策损失选类 & 行动或训练类别 & 类别支持程度与冲突强度 \\
    局部到整体 & 按任务规定的聚合规则转换 & 规则定义的整体结论 & 未保留的局部位置与数量 \\
    异构表示到共同空间 & 依据配对锚点、生成关系或可验证映射 & 映射覆盖的共同结构 & 尚未确定的语义对应 \\
    \bottomrule
  \end{tabularx}
  \caption{知识形式转换的条件与信息损失。“丢失”指转换结果不再保留的信息，“待补”指输入原本就不足以确定的信息；二者都不能从转换后的结果凭空恢复。}
  \label{tab:knowledge-fusion-form-conversions}
\end{table*}
\FloatBarrier

\section{融合知识的推断}
\label{sec:knowledge-fusion-inference}

现在把表~\ref{tab:knowledge-fusion-source-records}中的输出视为已对齐标签：同一行对应同一照片，各类别采用共同空间$Y$。在固定这些输入的条件下，我们比较三种证据组合办法：投票给每份有效输出相同贡献，加权改变贡献大小，概率推断则比较各候选真值产生当前输出的可能性。多个局部答案还可能受整体结构约束。这里比较的是计算与假设的差异，不是凭算法复杂程度评定质量；这些机制也可以组合。

\subsection{基于投票的知识融合}
\label{sec:knowledge-fusion-voting}

若暂时没有可用的能力参数，我们可以先问每个类别得到多少来源支持。\term{多数投票}（majority voting）把每份已对齐有效输出计为一票，再从共同集合$Y$中选择票数最大的类别。对样本$i$，需要先算出各候选$k$的票数，才能执行这个比较：
\begin{equation}
  v_i(k)
  =\sum_{s:o_{i,s}=1}\mathbb{I}\{\lambda_{i,s}=k\},
  \qquad
  \widehat y_i\in\argmax_{k\in Y}v_i(k).
  \label{eq:knowledge-fusion-majority-vote}
\end{equation}
求和只遍历$o_{i,s}=1$的来源，指示函数只给来源所报类别加一；没有有效输出的位置不读取$\lambda_{i,s}$。$\argmax$选出最高票的候选，可能不止一个。多类别的相对多数只要求票数最多，绝对多数还要求超过有效票的一半。若转向多标签任务，可以逐标签计票，但各项选择的组合仍要检查互斥、层级或序列条件；关系知识则须先对齐成同一有向命题，不能直接沿用实体类别票。

最高票相同或没有有效票时，式中的候选选择还不足以给出可用结果，需要另定接收规则。表~\ref{tab:knowledge-fusion-comparison}使用相对多数：唯一最高票即输出，平票或零票记为未决，没有另设最低票数。票数比例以有效票为分母，弃权和失败不计类别票；它描述来源赞同比例，不是校准后的真实正确概率。

\begin{table*}[htbp]
  \centering
  \small
  \renewcommand{\arraystretch}{1.15}
  \begin{tabularx}{\linewidth}{@{}ccccc>{\raggedright\setlength{\parindent}{0pt}\arraybackslash}X@{}}
    \toprule
    \textbf{样本} & \textbf{有效票} & \textbf{山茶/茶梅票} & \textbf{等权结果} & \textbf{加权结果} & \textbf{关键原因} \\
    \midrule
    $x_1$ & $4$ & $4/0$ & 山茶 & 山茶 & 全部来源一致 \\
    $x_2$ & $3$ & $2/1$ & 山茶 & 茶梅 & 模型权重超过两条规则之和 \\
    $x_3$ & $4$ & $3/1$ & 山茶 & 山茶 & 人工与规则共同支持山茶 \\
    $x_4$ & $1$ & $1/0$ & 山茶 & 山茶 & 只有模型作答 \\
    $x_5$ & $3$ & $2/1$ & 山茶 & 茶梅 & 两张山茶票共享模板 \\
    $x_6$ & $4$ & $2/2$ & 未决 & 茶梅 & 人工与模型共同支持茶梅 \\
    $x_7$ & $0$ & $0/0$ & 未决 & 未决 & 全部来源缺失或弃权 \\
    $x_8$ & $4$ & $2/2$ & 未决 & 山茶 & 人工与模型共同支持山茶 \\
    \bottomrule
  \end{tabularx}
  \caption{同一组候选的融合结果对照。加权列使用教学权重$w_{R_1}=w_{R_2}=0.25$、$w_H=1.2$、$w_M=0.6$，只用于复核运算与展示假设差异，不表示这些权重经过真实审计，也不证明加权方法普遍优于投票。}
  \label{tab:knowledge-fusion-comparison}
\end{table*}
\FloatBarrier

逐行复核表中的票数，可以得到平票、零票和唯一最高票三种结果。$x_5$还暴露了等权处理的代价：复制一条规则便把一张山茶票变成两张，却没有新增独立证据。投票便于作为基线，是因为无需拟合来源参数、计算可直接核对；它同时略去了能力、错误方向和依赖差异。新增来源是否改善正确性仍取决于这些差异及覆盖，不能从票数增加推出。

计算成本也要区分累计票数与选择类别。遍历$n\times m$稠密矩阵累计票数需$O(nm)$，若随后逐一比较$K$类，总时间为$O(nm+nK)$；只存有效输出时为$O(N_{\mathrm{obs}}+nK)$，其中$N_{\mathrm{obs}}=\sum_{i,s}o_{i,s}$。逐样本计算只需$O(K)$工作内存，保留全部票数需$O(nK)$。这些界只核算计票与选类，不包含语义对齐和相关来源去重。

\subsection{基于加权的知识融合}
\label{sec:knowledge-fusion-weighting}

如果已有审计或来源模型表明几份输出贡献不同，每列一票就没有利用这项知识。\term{加权融合}（weighted fusion）把有效输出乘以明确的贡献系数，再比较各类累计得分。固定前面的标签输入，对来源$s$在样本$i$支持类别$k$的贡献记为$w_{i,s,k}$，硬标签得分与选类规则为
\begin{equation}
  s_i(k)
  =\sum_{s:o_{i,s}=1}w_{i,s,k}
    \mathbb{I}\{\lambda_{i,s}=k\},
  \qquad
  \widehat y_i\in\argmax_{k\in Y}s_i(k).
  \label{eq:knowledge-fusion-weighted-vote}
\end{equation}
与等权计票相比，这里仍在$Y$中选择最高支持的类别，但支持单位改成了权重得分。$w_{i,s,k}$可以全局固定，也可以随类别、区域或时间变化，必须说明来自哪批审计、哪个来源模型或哪条设计规则，以及得分的量纲。选权重所用的数据属于开发依据，不能再充当最终质量评价。

硬标签只支持一个类别；若来源给出$p_s(k\mid x_i)$，则可以保留它分给各类的全部质量。我们改用同一照片、同一类别定义下的概率输出，并以非负来源权重$w_{i,s}$作线性汇总；此处$o_{i,s}=1$表示该来源给出了可用的完整类别分布。除以有效权重之和，是为了让结果仍在各类上和为一：
\begin{equation}
  \widetilde p_i(k)
  =\frac{\sum_{s:o_{i,s}=1}w_{i,s}p_s(k\mid x_i)}
  {\sum_{s:o_{i,s}=1}w_{i,s}}.
  \label{eq:knowledge-fusion-weighted-probability}
\end{equation}
先筛出给出完整可用分布的来源，再计算其加权概率，不对未返回分布的位置求值。各来源必须使用可比较的类别语义、概率尺度与校准含义，权重非负且分母为正；没有有效正权重时，应保留未决，而不是计算$0/0$。式~\eqref{eq:knowledge-fusion-weighted-probability}给出的是线性混合分布，不是把来源似然相乘的生成推断。任意模型分数不能直接当概率，分布和为一也只满足形式要求，不自动保证与真实频率校准。

回到表~\ref{tab:knowledge-fusion-comparison}的硬标签对照，我们预设$w_{R_1}=w_{R_2}=0.25,w_H=1.2,w_M=0.6$，让同模板规则合计贡献较低。$x_5$上山茶得$0.25+0.25=0.5$，茶梅得$1.2$，因而加权选择茶梅。表中其他行同样可以复算；权重只是教学配置，未从前四行的正确率拟合，结果反转不证明人工答案正确。

权重若直接取当前最大概率，就会把输出集中程度当作来源能力。为检查这个替代，设一张参考为茶梅的教学照片上，偏向山茶的$C$给出$p_C=(0.99,0.01)$，独立的$D$给出$p_D=(0.25,0.75)$，分量依次为山茶、茶梅。用最大概率$0.99,0.75$加权，两类未归一化得分为$1.1676,0.5724$，选出的山茶错误；假设独立审计支持权重$0.2,0.8$，得分变为$0.398,0.602$，选茶梅。两次计算固定分布和参考，只改变权重依据，说明集中输出也能集中在错误类别上。审计权重仍须核对代表性，概率本身仍须检查校准，不能由本例推导普遍收益。

把权重细化到类别和区域，需要对应的参考与覆盖记录，也增加更新维护成本。负权重把来源输出解释为反向证据，需要额外模型支持；它不适用于前面的概率线性混合。可以通过删除单个来源、合并相关来源、更新审计或扰动权重，检查类别选择对配置的敏感程度。加权表达了贡献不同，但尚未解释潜在真值如何产生观察，下一种方法将补上这个生成关系。

\subsection{基于概率模型与约束的知识融合}
\label{sec:knowledge-fusion-probabilistic-constrained}

当A报山茶、B报月季时，混淆参数能让我们比较哪一真实类别更可能产生这组观察；当序列各位置已有最高分类，规则还可能禁止它们组成当前序列。概率模型处理前一种问题，把真实标签设为未观测变量；结构约束处理后一种问题，限定或惩罚整体组合。两者合用时，可以先算局部概率支持，再在硬约束或软约束下选择一组标签，而不是把概率高与整体合法当作同一要求。

\subsubsection{潜在真实标签的概率推断}
\label{sec:knowledge-fusion-latent-label-inference}

面对已对齐但没有参考真值的标签矩阵，我们需要同时解释样本类别与来源错误。\term{Dawid--Skene模型}（Dawid--Skene model）把真实类别$Y_i$作为潜变量，也就是未直接观测而由模型推断的量，再用混淆矩阵描述各来源怎样报告它。设$\pi_k=\mathbb{P}(Y_i=k)$为类别先验，$\theta^{(s)}_{k,\ell}$沿用式~\eqref{eq:knowledge-fusion-confusion-matrix}，$S_i$是实际给样本$i$提供有效类别的来源集合。标准模型要求样本独立同分布、混淆行为在建模范围内稳定、给定$Y_i$后各来源独立。不另建缺失模型时，还需$S_i$由事前外生设计固定，或缺失机制对当前似然可忽略；若作答选择依赖未知类别或难度，事后看到$S_i$不等于可以忽略这种选择。固定这些条件后，先验乘各来源似然给出候选真值与观察的联合概率，再对全部$K$类求和，才得到只含已见标签的概率：
\begin{equation}
  p(Y_i=k,\bm{\lambda}_i)
  =\pi_k\prod_{s\in S_i}\theta^{(s)}_{k,\lambda_{i,s}},
  \qquad
  p(\bm{\lambda}_i)
  =\sum_{k=1}^{K}\pi_k
    \prod_{s\in S_i}\theta^{(s)}_{k,\lambda_{i,s}}.
  \label{eq:knowledge-fusion-dawid-skene}
\end{equation}
第一项将类别先验与该类别下的观察可能性合并，第二项对真值求和后给出当前观察的总概率，可用于比较候选参数对数据的解释。在先验非负且和为一、混淆矩阵逐行归一化的参数集合中，我们要最大化各样本观测概率乘积，通常等价地最大化对数似然。求和位于对数内部，不能直接按已知真值计数；Dawid与Skene采用EM算法联合估计观察者误差率与潜在类别\scite{dawid1979observer}

这里沿用前卷期望最大化（expectation-maximization, EM）的思路：先按当前参数分配未知类别的概率，再用这些概率形成软计数更新参数。E步需要的量是责任度（responsibility）$q_{i,k}^{(t)}$，即第$t$轮模型看到来源输出后认为样本$i$属于$k$的后验概率。用候选类别的联合概率除以所有类别联合概率之和，就得到可供下一步计数的权重：
\begin{equation}
  q_{i,k}^{(t)}
  =
  \frac{\pi_k^{(t)}
    \prod_{s\in S_i}
    \left(\theta^{(s)}_{k,\lambda_{i,s}}\right)^{(t)}}
  {\sum_{r=1}^{K}\pi_r^{(t)}
    \prod_{s\in S_i}
    \left(\theta^{(s)}_{r,\lambda_{i,s}}\right)^{(t)}}.
  \label{eq:knowledge-fusion-dawid-skene-e-step}
\end{equation}
分母须为正；它比较的仍是共同集合$Y$中的全部类别，而不是只比较被投过票的类别。若$S_i$为空，空乘积为一，在当前可忽略缺失假设下后验退回先验，没有从本行获得类别证据。M步固定$q^{(t)}$，最大化完整数据对数似然的后验期望：每个样本以$q_{i,k}^{(t)}$贡献给类别$k$，有输出$\ell$的样本再贡献给来源矩阵的第$k,\ell$项。在先验和每行混淆均归一化的约束下，加权计数除以相应总量得到
\begin{equation}
  \pi_k^{(t+1)}
  =\frac{1}{n}\sum_{i=1}^{n}q_{i,k}^{(t)},
  \qquad
  \left(\theta^{(s)}_{k,\ell}\right)^{(t+1)}
  =
  \frac{\sum_{i:s\in S_i}
    q_{i,k}^{(t)}\mathbb{I}\{\lambda_{i,s}=\ell\}}
  {\sum_{i:s\in S_i}q_{i,k}^{(t)}}.
  \label{eq:knowledge-fusion-dawid-skene-m-step}
\end{equation}
先验更新按全部$n$个样本平均类别责任度，混淆更新则只使用来源$s$实际作答的样本；后一个分母是这些样本贡献给真值$k$的软计数。它为零时，这一行未得到有效权重，应保留未估计状态或加入明确先验，不能计算$0/0$。多数票、少量审计或接近对角的矩阵可以初始化，但分别带入不同方向假设。迭代应同时设最大次数与似然或参数变化阈值；共享错误、选择性覆盖或未观测类别--来源组合仍要另查，更新式不会消除这些模型与证据缺口。

为了复算一次更新，我们另取三个始终作答的来源，保持山茶、茶梅二分类，初始对称正确率为$0.8,0.7,0.6$，先验各为$0.5$。三个教学样本的观察为$(\text{山茶},\text{山茶},\text{山茶})$、$(\text{山茶},\text{茶梅},\text{茶梅})$、$(\text{茶梅},\text{茶梅},\text{山茶})$。第一行在山茶与茶梅下的权重为$0.168$与$0.012$，分别来自$0.5\times0.8\times0.7\times0.6$和$0.5\times0.2\times0.3\times0.4$；第二行两类权重为$0.048,0.042$，第三行为$0.018,0.112$。逐行归一化，山茶责任度为$0.9333,0.5333,0.1385$，茶梅为各自补数。第二行虽有两票茶梅，当前能力参数却给山茶略高后验，展示了来源似然与票数的区别。

用未舍入责任度更新，山茶与茶梅的软计数总量约为$1.6051,1.3949$，除以三个样本得到先验$(0.5350,0.4650)$。来源1在前两行报山茶，这两行分给真值山茶的权重合计为$0.9333+0.5333$，除以山茶总软计数得$0.9137$；同样，两行分给真值茶梅的权重合计约$0.5334$，除以茶梅总软计数得$0.3824$。这些分数不是参考计数，而是当前后验下的期望计数。记第1次M步后的矩阵为$\bm{\Theta}_s^{(1)}=\left(\left(\theta^{(s)}_{k,\ell}\right)^{(1)}\right)_{k,\ell}$，按行是真值山茶/茶梅、列是输出山茶/茶梅排列，得到
\[
  \begin{aligned}
    \bm{\Theta}_1^{(1)}
    &=
    \begin{pmatrix}
      0.9137 & 0.0863 \\
      0.3824 & 0.6176
    \end{pmatrix},\\
    \bm{\Theta}_2^{(1)}
    &=
    \begin{pmatrix}
      0.5815 & 0.4185 \\
      0.0478 & 0.9522
    \end{pmatrix},\\
    \bm{\Theta}_3^{(1)}
    &=
    \begin{pmatrix}
      0.6677 & 0.3323 \\
      0.6654 & 0.3346
    \end{pmatrix}.
  \end{aligned}
\]
逐行求和均为一，先验与三个来源的完整混淆矩阵都已更新，而不只是更新一个总体正确率。初值设为对称误差，M步没有强制继续对称，所以两行正确率可以不同。这一轮未加平滑，三个样本仅用于复核运算，不足以稳定恢复来源能力。

图~\ref{fig:knowledge-fusion-latent-label-em}将这轮计算接成循环：标签矩阵固定，E步改责任度，M步改参数，再检查是否停止。若在第$t$轮M步后停止，最终先验和混淆参数已是$\pi_k^{(t+1)}$与$\left(\theta^{(s)}_{k,\ell}\right)^{(t+1)}$，而$q_{i,k}^{(t)}$仍对应更新前的参数；交付前必须用最终参数按式~\eqref{eq:knowledge-fusion-dawid-skene-e-step}重算$q_{i,k}^{(t+1)}$。例如只做上述一轮M步就结束，用未舍入参数刷新后，三个样本的山茶后验约为$0.9711,0.5455,0.0662$，不同于更新前的$0.9333,0.5333,0.1385$。图中“后验分布与来源参数”的出口包含这次刷新，返回相匹配的参数与后验，并记录停止原因。

精确E步与M步不会降低当前模型的观测似然，但目标一般非凸，数值趋稳不保证全局最优，也可能发生在局部最优或鞍点附近。达到最大迭代次数是资源上限触发的停止，不代表已经收敛；它与图中的稳定判断是不同停止原因。即使更新稳定，也只回答数值迭代是否趋稳，不回答参数是否可识别，更不证明推断类别符合现实真值。

\begin{figure}[!t]
  \centering
  \begin{tikzpicture}[
  x=1mm,y=1mm,
  every node/.style={font=\figurefont\small,text=BookInk},
  data/.style={book node,draw=BookBlue,fill=FigureInput!10,line width=0.9pt,minimum width=42mm,minimum height=11mm,inner sep=2pt},
  state/.style={book node,draw=BookTeal,fill=FigureModel!10,line width=0.9pt,minimum width=45mm,minimum height=11mm,inner sep=2pt},
  check/.style={book node,draw=BookAmber,fill=BookAmber!10,line width=0.9pt,minimum width=34mm,minimum height=10mm,inner sep=2pt},
  output/.style={book node,draw=BookInk,fill=BookPaper,line width=0.9pt,minimum width=39mm,minimum height=10mm,inner sep=2pt}
]
  \node[data] (matrix) at (0,34) {固定观测$\bm{\Lambda},\bm{O}$};
  \node[state] (init) at (0,17) {初始化$\bm{\pi}^{(0)},\bm{\Theta}^{(0)}$};
  \node[state] (estep) at (0,0) {E步：计算$q_{i,k}^{(t)}$};
  \node[state] (mstep) at (0,-17) {M步：更新$\bm{\pi}^{(t+1)},\bm{\Theta}^{(t+1)}$};
  \node[check] (stop) at (0,-34) {似然或模型参数已稳定？};
  \node[output] (result) at (0,-51) {后验分布与来源模型参数};

  \draw[book arrow,knowledge arrow,draw=BookBlue,line width=1pt] (matrix) -- (init);
  \draw[book arrow,knowledge arrow,draw=BookTeal,line width=1pt] (init) -- (estep);
  \draw[book arrow,knowledge arrow,draw=BookTeal,line width=1pt] (estep) -- node[right,font=\figurefont\scriptsize] {软类别计数} (mstep);
  \draw[book arrow,knowledge arrow,draw=BookTeal,line width=1pt] (mstep) -- (stop);
  \draw[book arrow,knowledge arrow,draw=BookInk,line width=1pt] (stop) -- node[right,font=\figurefont\scriptsize] {是} (result);
  \draw[book arrow,knowledge arrow,draw=BookBlue] (matrix.east) -- ++(16,0) |- node[pos=0.22,right,font=\figurefont\scriptsize,text=BookBlue] {每轮读取} (estep.east);
  \draw[book arrow,knowledge arrow,draw=BookAmber,dashed] (stop.west) -- ++(-16,0) |- node[pos=0.18,left,font=\figurefont\scriptsize,text=BookAmber,align=right] {否：携带\\新模型参数} (estep.west);

  \node[anchor=west,font=\figurefont\scriptsize,text=BookMuted,align=left] at (32,-42) {分别检查：\\可识别性\\数值收敛\\答案正确性};
\end{tikzpicture}

  \caption{来源参数与潜在标签的交替估计。观测标签矩阵保持不变；E步计算当前参数的类别后验，M步用软计数更新参数。稳定判断后的输出框包含用最终参数重算后验的刷新，再交付相匹配的后验与来源参数；图未展开的最大次数出口也执行这次刷新，但资源停止不表示收敛。实线表示数据与计算流，虚线表示新参数回流。}
  \label{fig:knowledge-fusion-latent-label-em}
\end{figure}

回到第~\ref{sec:knowledge-fusion-relational-reliability}节的“全部正确或全部报反”，Dawid--Skene也保留这种类别置换：同时交换先验元素、所有混淆矩阵的行与潜在类别名称，观测分布$p(\bm{\Lambda})$不变。参考锚点或方向假设负责把潜在类别命名为真实物种，重叠不足与退化矩阵还可能留下其他歧义。来源实际相关却按独立相乘，是另一种问题：它破坏生成假设，即使拟合稳定也会重复使用证据。

若希望知道类别置换之外能否唯一恢复参数，我们还需检查不同类别在来源输出中留下了多少可区分信息。下面对一类来源模型给出充分条件，不把它当作所有融合方法的必要条件。将先验与来源参数统记为$\bm{\theta}$，严格可识别要求相同的观测分布只对应同一参数：
\[
  p_{\bm{\theta}}(\bm{\Lambda})
  =p_{\bm{\theta}'}(\bm{\Lambda})
  \quad\Longrightarrow\quad
  \bm{\theta}=\bm{\theta}',
\]
若预先把类别名称置换视为等价，结论相应放宽到这类变换。还须每类先验$\pi_k>0$：先验为零的类别不参与观测分布，无论把它的来源输出概率改成什么，观察都不变，因此无法恢复该行参数。

% 保留矩阵定义的引导段与首个公式，避免前一幅浮动图插入句中。
\Needspace{45mm}
我们把来源分成三个非空组$G_1,G_2,G_3$，以容许组内共享信息。给定$Y_i$后三组联合输出须独立，组内则可以依赖。下面的识别结论针对完整的三组联合输出分布：实际记录须在同一样本上联合观测三组，或另有已成立的观测与恢复条件，能从缺失记录识别这个完整分布。仅有各组边缘分布、覆盖连通或可忽略缺失，还不能据此保证完整联合信息可恢复。

当前比较单位从单列输出改为组内所有输出构成的向量$\bm{\lambda}_{i,G_j}$。这些联合状态组成向量集合$\bm{\Omega}_j$，也可包含模型明确保留的无票状态；每个状态$\bm{u}$按组内来源顺序列出输出。要衡量第$j$组能否区分类别，先定义$\bm{M}_j\in\mathbb{R}^{K\times|\bm{\Omega}_j|}$：
\[
  [\bm{M}_j]_{k,\bm{u}}
  =\mathbb{P}(\bm{\lambda}_{i,G_j}=\bm{u}\mid Y_i=k),
  \qquad
  k\in Y,\ \bm{u}\in\bm{\Omega}_j.
\]
每行是一个潜在类别下的组输出分布；行相同表示这组输出无法区分相应类别。行Kruskal秩（row Kruskal rank）$\operatorname{krank}(\bm{M}_j)$是任取$r$行都线性无关的最大$r$，用来检查类别分布之间的线性区别。固定三组输出后，每个类别贡献其先验乘三个组条件概率，再对类别求和得到观测概率；这就是以类别为混合成分的三因子分解。下面的秩和门槛为这种分解提供除成分置换外的唯一性条件：
\begin{equation}
  \operatorname{krank}(\bm{M}_1)
  +\operatorname{krank}(\bm{M}_2)
  +\operatorname{krank}(\bm{M}_3)
  \geq 2K+2,
  \label{eq:knowledge-fusion-kruskal-condition}
\end{equation}
在完整三组联合输出分布可由观测确定或按已成立机制恢复、每类先验为正、三组条件独立且满足式~\eqref{eq:knowledge-fusion-kruskal-condition}时，当前参数点的类别先验与三个组条件矩阵除类别置换外可识别。这里引用三因子分解的唯一性条件；概率行归一化消除了分解中的缩放歧义，但仍须另用锚点确定类别名称。这项结论不使仅有各组边缘的记录自动具备联合信息，也没有给有限样本估计误差保证\scite{allman2009identifiability}

直接核对实际矩阵的秩，得到的是当前参数点的结论。若只知道每组联合空间有多大，还可以问这些维度是否容许足够高的秩；每个矩阵的最大行Kruskal秩至多为$\min\{K,|\bm{\Omega}_j|\}$。当维度满足
\[
  \sum_{j=1}^{3}\min\{K,|\bm{\Omega}_j|\}\geq2K+2,
\]
在一般位置上，三个条件概率矩阵各自达到允许的最大行Kruskal秩，前面的逐点门槛因而一般成立。“一般位置”排除了代数退化构成的零测集合，不等于每个参数点都成立，也不排除有限样本不稳定。实际模型有额外参数约束时，必须检查可产生的矩阵；仅数输出空间大小可能高估区分能力。

复制来源没有新增信息，但是否使已有区分能力下降，要看实际行分布。设一个二分类来源对称正确率为$0.8$，组内第二个来源逐项复制它；以$0,1$编码两类，联合输出依次为$00,01,10,11$。不一致的两列概率为零，该组条件矩阵为
\[
  \bm{M}
  =\begin{pmatrix}
    0.8 & 0 & 0 & 0.2\\
    0.2 & 0 & 0 & 0.8
  \end{pmatrix}.
\]
第1、第4列的二阶子式为$0.8^2-0.2^2=0.6$，所以两行独立，行Kruskal秩仍为$2$。若三个这样的组给定真值后独立，秩和为$6=2K+2$，满足二分类门槛；复制没有增秩，也没有降低本例已有的秩。把正确率换成$0.5$，两行相同，秩才降到$1$；与类别无关的恒定输出同样只有秩$1$。组间依赖则破坏三因子分解前提，不是一个单纯的秩问题。来源数量、组空间大小与EM收敛都不能替代这些检查。

前面的乘积仅使用实际作答的来源；如果成功运行却没有类别票的倾向本身需要建模，或已知某些输出共享影响，就要把这些状态与关系也写入分布。\term{标签模型}（label model）是描述弱来源怎样产生输出的生成模型家族，不是Dawid--Skene的别名。在这一模型接口中，令标注函数输出$\lambda_{i,s}\in Y\cup\{\bot\}$，并将$\bot$约定为成功执行后的无类别票符号。与表~\ref{tab:knowledge-fusion-source-records}中专指主动弃权的符号相比，这里扩大了计算编码的范围：主动弃权和规则未触发都可映为$\bot$，但原记录的前者动作仍为弃权，后者动作仍为空、触发结果仍为未触发，内容均为未知。这个合并编码不能区分两种无票原因；需要区分时须加入原动作与触发结果，或在模型输出域另设状态。用$R$个函数$\phi_r$描述类别与来源输出的特定关系，再以参数$\theta_r$控制其强度，一个通用对数线性形式为
\begin{equation}
  p_{\bm{\theta}}(Y_i,\bm{\lambda}_i)
  =\frac{1}{Z_{\bm{\theta}}}
  \exp\left(
    \sum_{r=1}^{R}\theta_r
    \phi_r(Y_i,\bm{\lambda}_i)
  \right).
  \label{eq:knowledge-fusion-label-model}
\end{equation}
指数函数给每个$(Y_i,\bm{\lambda}_i)$配置正权重，$Z_{\bm{\theta}}$是这些权重在全部类别和输出配置上的和，除以它才成为概率分布。例如$\phi_r=\mathbb{I}\{\lambda_{i,s}=Y_i\}$记录来源是否报对，相应参数为正时提高这类配置的相对权重。因子还可以计量类别先验、偏好、覆盖、无票状态或来源对的相关行为；最简版本只保留先验与条件独立准确性因子，依赖版本再加入关系。数据编程（Data Programming）将规则、知识库和模型写成可以不投票的标注函数，再利用重叠与冲突拟合生成参数；其中规则条件不匹配时的无票输出，也不要求在原记录中伪造主动弃权\scite{ratner2016data-programming}

用式~\eqref{eq:knowledge-fusion-label-model}无标签拟合时，仍要对真值求和。单个样本的观测概率为$p_{\bm{\theta}}(\bm{\lambda}_i)=\sum_kp_{\bm{\theta}}(Y_i=k,\bm{\lambda}_i)$。选定因子结构与允许参数范围后，要最大化观测的对数似然总和$\sum_i\log p_{\bm{\theta}}(\bm{\lambda}_i)$。拟合后用$p_{\bm{\theta}}(Y_i=k,\bm{\lambda}_i)/p_{\bm{\theta}}(\bm{\lambda}_i)$计算类别后验，分母须为正。这个目标可回接刚才的EM框架：E步计算当前参数的类别后验，M步最大化完整数据对数似然的后验期望。不过，对数线性形式的M步包含依赖参数的$\log Z_{\bm{\theta}}$，须按所选因子结构确定数值优化或近似方法，不能套用式~\eqref{eq:knowledge-fusion-dawid-skene-m-step}的混淆行计数解。本段止于模型形式、拟合目标与求解框架的对应，不给通用因子模型的解析更新；直接枚举归一化需要遍历$K(K+1)^m$种类别与输出配置，加入依赖后能否利用结构降低成本，也须另查。

上述观测目标假设固定参数下样本独立、来源行为在建模范围内稳定；它只读取输出矩阵，不自动从原始$\bm{x}_i$识别困难照片。质量随样本或时间变化时，需要分组或加入条件因子。未请求、未处理、处理中和执行失败都不是这里的$\bot$；若这些状态进入模型，必须另写观测设计或生成机制，不能借成功执行后的无票因子一并解释。

% 段末文献需要与本段一同检查可用高度。
\Needspace{70mm}
固定同一类别任务比较两者，Dawid--Skene通常给每个来源完整混淆矩阵，程序化弱监督的标签模型则常用较少的准确性因子，并显式描述无票状态或依赖。两者都推断潜在类别，但改变的是来源输出的参数化与生成假设，不是“标签模型”替换了一个名称。Snorkel将这类生成模型得到的概率标签用于训练下游判别模型，后者再学习从原始输入预测类别\scite{ratner2017}

保留完整混淆能够区分每种错误方向，也需要更多分类别重叠证据；较粗的准确性因子减少参数，代价是加入对称误差或共享结构等限制。选择哪种描述，要看当前分歧是否需要这些方向差异，以及数据能否支持估计。

增加因子也会增加无法区分参数的风险。除类别置换外，对数线性表示可能冗余，使不同参数给同一分布；条件独立版本需要足够重叠与有效独立来源，依赖版本还需因子结构正确且证据能区分其参数。Data Programming的恢复结论只适用于论文规定的模型和条件，不意味着任意规则都能恢复真值。高度相关、类别方向未知或参数超过有效证据时，应补锚点、简化或合并来源，必要时保留未决；后验收敛不是验证证据。

\subsubsection{结构约束下的联合推断}
\label{sec:knowledge-fusion-structured-inference}

设一句话每个词元都已有多源融合得分，逐位置选最高类别却可能得到不合法实体跨度。BIO（beginning--inside--outside）标注用$B$表示实体开始、$I$表示同类实体内部、$O$表示实体外部，$I$不能在没有相应开始的位置接续实体。类别层级和互斥规则也会限制多个答案的组合。\term{联合推断}（joint inference）就是一起选择这些相互约束的标签，而不是分别选完再假定整体有效。

与前面的规则获取相比，本节的输入已经改变。规则获取由已有条件生成候选知识，本节已经有多源局部得分，要从候选组合选出整体赋值。相同业务规则可以参与两处过程，但前者评价生成的候选，后者评价整体可行性和得分，不能因都使用约束就混成同一步。

要在局部支持与结构要求间取舍，需要明确哪些组合根本不允许，哪些只应受罚。设待选标签为$z_i\in Y$、局部得分为$a_i(z_i)$；任务硬约束定义可行向量集合$\bm{Z}_{\mathrm{feas}}$，软约束$r$的违反程度为$\rho_r(\bm{z})\geq0$，罚值权重$\eta_r\geq0$须使代价与得分可比较。在允许集合中最大化支持总分并减去违反代价，得到
\begin{equation}
  \widehat{\bm{z}}
  \in\argmax_{\bm{z}\in\bm{Z}_{\mathrm{feas}}}
  \left\{
    \sum_i a_i(z_i)
    -\sum_r\eta_r\rho_r(\bm{z})
  \right\}.
  \label{eq:knowledge-fusion-structured-objective}
\end{equation}
硬约束先排除集合外候选，软约束则在剩余候选间降低违反规则的得分；罚号为负，是因为这里最大化目标。局部分布为$q_{i,k}$时可以取$\log q_{i,k}$作为得分，零概率对应$-\infty$；若没有整体概率分解依据，这仍是评分目标，而非自动得到联合后验。未校准加权分数同样只表达当前评分体系的偏好。

我们用三个词元枚举一次联合选择。共同标签为$O,B\text{-ORG},I\text{-ORG}$，后两者表示组织实体的开始与内部；第三个位置只允许$O$。表内$-\infty$排除对应位置标签，其余数值是人为设定的局部支持分：
\[
  \begin{array}{c|ccc}
    & O & B\text{-ORG} & I\text{-ORG} \\
    \hline
    1 & 1.2 & 0.8 & -\infty \\
    2 & 0.5 & 0.6 & 1.4 \\
    3 & 0.7 & -\infty & -\infty
  \end{array}.
\]
逐位置最大化选出$(O,I\text{-ORG},O)$，总分$3.3$，却违反了$I\text{-ORG}$不能接在$O$后的规则。有限分数的合法组合只有五个：$(O,O,O)$为$2.4$，$(O,B\text{-ORG},O)$为$2.5$，$(B\text{-ORG},O,O)$为$2.0$，$(B\text{-ORG},B\text{-ORG},O)$为$2.1$，$(B\text{-ORG},I\text{-ORG},O)$为$2.9$。硬约束选择最后一个，它改变第一个位置以保留第二个位置的强证据，并非仅删除冲突标签。

保持局部分数不变，只把这条转移规则改成一次违反记$1$的软代价，非法序列便以$3.3-\eta$参加比较，其余五个候选罚值为零。$\eta=0.2$时它得$3.1$，仍胜过合法最优$2.9$；$\eta=0.5$时只得$2.8$，转而选择$(B\text{-ORG},I\text{-ORG},O)$。硬约束改变可选集合，软约束改变得分，两者在相同输入上有不同作用。互斥分类也可以保留得分损失较小的一项；若互斥规则可能不适用，则需明确是否改为允许违反的软约束。

选择合法组合仍以约束正确且适用为前提。约束可能过时或互相冲突，求解失败也可能是确无可行赋值、当前样本被施加了不适用规则，或求解器在资源限额内未找到解。删除冲突应保留原来源，放松约束应记录违反程度，近似解要说明尚未保证可行性还是最优性。链式序列可用动态规划；更一般的离散耦合可能需要整数规划、消息传递或近似搜索，不能由模型存在推断求解成本可承受。

\section{融合结果的表达与使用}
\label{sec:knowledge-fusion-result-use}

假设融合给山茶分配$0.62$的支持，一项训练任务可以保留这份分布，高风险使用却可能不接收任何类别。差异来自用途，而非同一模型同时有两种确信程度。我们要把票数、得分、后验或整体赋值这些推断结果，与选类、训练、保留候选或弃权这些使用动作分开，并保存来源摘要、对齐映射、融合配置和结果版本。

\subsection{单一标签与概率分布的使用}
\label{sec:knowledge-fusion-label-probability-use}

如果接收端必须选一个类别，后验最高的类别是否合适，还取决于错误代价。错把茶梅报成山茶与漏掉某个需要复核的类别，可能有不同损失，单看最大概率没有表达这种取舍。设$q_i(k)$为融合类别分布，$A$为用途允许的行动集合，$L(a,k)$为真值$k$下行动$a$的损失；我们对未知真值按$q_i$加权平均损失，再在允许行动中选最小者，得到贝叶斯决策：
\begin{equation}
  \widehat a_i
  \in\argmin_{a\in A}
  \sum_{k=1}^{K}L(a,k)q_i(k).
  \label{eq:knowledge-fusion-bayes-decision}
\end{equation}
这一求和回答每个行动在当前后验下平均要付出多少损失，$\argmin$的方向是损失越小越好。行动集合就是类别且采用零一损失时，期望损失为$1-q_i(a)$，因此才等价于最大后验。另设一个二分类教学后验$q_i=(0.62,0.38)$，分量为山茶、茶梅；若报错山茶的损失为$2$、报错茶梅为$1$、报对为零，选山茶的期望损失为$2\times0.38=0.76$，选茶梅为$1\times0.62=0.62$，决策便选茶梅。同一后验在零一损失下选山茶，在这组损失下选茶梅，改变的是用途代价，不是模型确信程度。

接收端若是分类模型训练，目的不再是立即执行一个行动，而是让预测分布接近融合分布。可以用软标签交叉熵$-\sum_kq_i(k)\log p_{\bm{w}}(k\mid\bm{x}_i)$，其中$p_{\bm{w}}$是待训练模型的类别概率，$q_i$让各类支持都进入损失。只保留$\argmax_kq_i(k)$则把剩余质量舍去，接口与存储更简单，却隐藏了冲突强度。由未校准分数或错设来源模型得到的$q_i$只表示当前模型支持，软训练不会自动把它变成真实频率。

分布与使用规则都要保留版本。新增来源、更新混淆或修订语义映射会改变推断，修改阈值或损失矩阵则只改变使用决策；审计时需要原分布与两类配置才能定位变化。训练接口只接收硬标签时，可以在接口处选类，仍保留上游概率记录供复核。

\subsection{候选集合与弃权的使用}
\label{sec:knowledge-fusion-set-abstention-use}

如果支持分布还不能区分两个类别，立即选一个会丢掉现有歧义，可以先输出候选集合$C_i\subseteq Y$。一种构造按$q_i(k)$降序加入类别，直到累计质量达到预设$\tau\in(0,1]$，相同概率的排序也要有固定规则；另一种保留所有$q_i(k)\geq\gamma$的类别，$\gamma\in[0,1]$是逐类接收阈值。前者在有效分布下会得到非空集合，后者可能为空；单元素集合相当于硬标签，过大集合则少有区分作用。累计质量或逐类概率达到阈值不等于覆盖真实类别，覆盖保证还需独立校准和相应假设。

保留候选后，使用端仍可能要求更确定的判断。综合证据不能形成当前用途可接收的结论时，结果记为“未决”；使用端主动不输出决策的动作沿用\term{弃权}的约定。它与输入来源运行后弃权处于不同环节，应分别记录。系统可以因$\max_kq_i(k)<\eta$、有效独立来源不足或对齐未确认而记未决，再请求复核、追加信息或由使用端弃权；这些接收条件需预先明确。弃权不是负类，其对象仍属于目标总体，覆盖与质量统计不能永久略去它们。

图~\ref{fig:knowledge-fusion-output-forms}固定教学分布$q=(0.62,0.28,0.10)$，分量依次为山茶、茶梅、月季，再只改变使用规则。最大后验选山茶，软训练保留三项；累计阈值$\tau=0.90$时，$0.62+0.28=0.90$，候选为$\{\text{山茶},\text{茶梅}\}$。若高风险用途要求最大支持至少$0.70$，当前$0.62$不足，记为未决并由使用端弃权。四个分支是不同用途的对照，不是一次使用中同时执行互斥行动。

\begin{figure*}[htbp]
  \centering
  \begin{tikzpicture}[
  x=1mm,y=1mm,
  every node/.style={font=\figurefont\small,text=BookInk},
  evidence/.style={book node,draw=BookInk,fill=BookPaper,line width=1pt,minimum width=61mm,minimum height=12mm,inner sep=2pt},
  result/.style={book node,minimum width=33mm,minimum height=18mm,inner sep=2pt,text width=30mm,align=center}
]
  \node[evidence] (q) at (75,19) {教学支持分布（山茶、茶梅、月季）\\$q=(0.62,\ 0.28,\ 0.10)$};

  \node[result,draw=BookBlue,fill=FigureInput!9,line width=0.9pt] (hard) at (12,-14) {硬标签\\山茶};
  \node[result,draw=BookTeal,fill=FigureModel!9,line width=0.9pt] (soft) at (54,-14) {概率目标\\$(0.62,0.28,0.10)$};
  \node[result,draw=BookAmber,fill=BookAmber!9,line width=0.9pt,dashed] (set) at (96,-14) {候选集合\\$\{\text{山茶},\text{茶梅}\}$};
  \node[result,draw=BookMuted,fill=BookPaper,line width=0.9pt,dotted] (abstain) at (138,-14) {暂缓使用\\弃权};

  \draw[book arrow,knowledge arrow,draw=BookBlue,line width=0.9pt] (q.south west) -- ++(0,-4) -| node[pos=0.74,left,font=\figurefont\scriptsize,text=BookBlue] {$\argmax$} (hard.north);
  \draw[book arrow,knowledge arrow,draw=BookTeal,line width=0.9pt] ([xshift=15mm]q.south west) -- ++(0,-6) -| node[pos=0.76,left,font=\figurefont\scriptsize,text=BookTeal] {完整保留} (soft.north);
  \draw[book arrow,knowledge arrow,draw=BookAmber,dashed,line width=0.9pt] ([xshift=-15mm]q.south east) -- ++(0,-6) -| node[pos=0.76,right,font=\figurefont\scriptsize,text=BookAmber] {累计权重$\geq0.90$} (set.north);
  \draw[book arrow,knowledge arrow,draw=BookMuted,dotted,line width=0.9pt] (q.south east) -- ++(0,-4) -| node[pos=0.74,right,font=\figurefont\scriptsize,text=BookMuted] {$\max q<0.70$} (abstain.north);

  \node[font=\figurefont\scriptsize,text=BookMuted] at (12,-31) {离散决策};
  \node[font=\figurefont\scriptsize,text=BookMuted] at (54,-31) {软标签训练};
  \node[font=\figurefont\scriptsize,text=BookMuted] at (96,-31) {部分排除};
  \node[font=\figurefont\scriptsize,text=BookMuted] at (138,-31) {触发复核};
\end{tikzpicture}

  \caption{同一融合证据的不同输出方式。三类别支持分布为人工教学设定；四个分支分别采用最大后验、完整概率、累计概率权重阈值和最低接收阈值，展示硬化、保留不确定性、保留候选与暂缓使用的差异，不赋予该分布额外校准或覆盖保证。}
  \label{fig:knowledge-fusion-output-forms}
\end{figure*}
\FloatBarrier

集合用于训练还需区分“模型筛出的候选”与“已确认真值在其中”的知识。假设照片已确认属于山茶或茶梅，但无法再细分，完整类别表还包括月季；将这个已确认集合记为$S(\bm{x}_i)$，不同于前面按分数构造的$C_i$。当前监督要求提高集合内总概率，而没有指定山茶和茶梅各占多少。对参数$\bm{w}$的分类模型输出$p_{\bm{w}}(k\mid\bm{x}_i)$，先在确认集合内求和，再以负对数惩罚总质量过小，得到集合约束损失：
\begin{equation}
  \ell_{\mathrm{set}}(\bm{w};i)
  =-\log\sum_{k\in S(\bm{x}_i)}
    p_{\bm{w}}(k\mid\bm{x}_i).
  \label{eq:knowledge-fusion-set-loss}
\end{equation}
集合内概率越大，负对数损失越小；总质量为零则损失发散。若三类概率为$(0.6,0.3,0.1)$，本例为$-\log(0.6+0.3)=-\log0.9\approx0.10536$。降低这项损失要求减少集合外月季的质量，却不区分集合内两类，也不要求各占一半。把两类分别当确定正类，或指定各一半软标签，都改变了原来的集合事件监督。

式~\eqref{eq:knowledge-fusion-set-loss}只使用已确认的集合事件，并没有描述候选集合怎样生成，不能视为任意候选机制的完整似然。模型筛出的$C_i$可能漏掉真值，直接当确定约束会迫使模型压低真类概率；这类记录应保留候选状态，按生成条件另定接收、降权或复核办法。

阈值还会改变评价对象。固定模型分数提高最低接收阈值，已接收集合只会缩小或不变，但其中错误比例不保证逐组下降；困难区域若被大量排除，总体准确率改善可能伴随服务范围收缩。因此要同时报告总体与分组覆盖、已接收错误率和集合大小，并把阈值选择数据与最终独立评价分开。

用一组人工计数固定这个比较：简单区与困难区各$100$个样本，阈值$0.60$时分别接收$90,60$个，错$9,18$个，总体已接收错误率为$27/150=18\%$。阈值改为$0.80$后接收$70,20$个，错$4,5$个，总体错误率为$9/90=10\%$。错误率下降了$8$个百分点，但总体覆盖也从$150/200=75\%$降到$90/200=45\%$，困难区覆盖从$60\%$降到$20\%$，已接收错误率为$5/20=25\%$。本例两区错误率都下降，不是反转例子；它说明只报总体错误改善会遗漏困难区的大幅覆盖损失。数字均为教学设定。

\subsection{异构知识的联合使用}
\label{sec:knowledge-fusion-heterogeneous-use}

一张花卉照片的类别、两张视图来自同一对象的关系，以及遮蔽花瓣的恢复内容，不能在同一类别表内计票。它们要求模型预测或保留不同对象。\term{异构知识的联合使用}（joint use of heterogeneous knowledge）就是保留各类监督对象，让它们通过共享参数、多个损失或结构约束共同影响训练；要组织的是联合训练安排，而非强制产生一个融合标签。

要把这些要求交给同一模型，先为每类知识定义与其内容匹配的损失，再确定哪些训练单元具有该知识。设本次训练有$n$个单元、$R$种知识，共享参数为$\bm{w}$，$\ell_r(\bm{w};i)$是知识类型$r$的单项损失，$m_{i,r}\in\{0,1\}$表示单元$i$是否适用，$\alpha_r\geq0$控制该项贡献。这里一个单元可以记录单张图像，也可以包含配对视图及其索引，不把关系误作单图类别。我们在各类有效单元上平均损失，再加权形成待最小化的经验目标：
\begin{equation}
  \mathcal{L}(\bm{w})
  =
  \sum_{r=1}^{R}
  \alpha_r
  \frac{\sum_{i:m_{i,r}=1}\ell_r(\bm{w};i)}
  {\max\left(1,\sum_{i=1}^{n}m_{i,r}\right)}.
  \label{eq:knowledge-fusion-heterogeneous-objective}
\end{equation}
先按$m_{i,r}=1$筛出具有相应知识的单元，再计算$\ell_r(\bm{w};i)$；缺少配对、类别或恢复内容时，不求值相应损失，也不填零目标。某类有效数为零时，分子是空和，按零处理，$\max(1,\cdot)$防止零分母，整项贡献为零，但该类知识仍属缺失；这不是用零标签监督模型。逐项平均消除了有效数的直接倍数影响，却未统一不同损失尺度，$\alpha_r$还需结合目标重要性、梯度尺度、采样频率与成本设定。

沿用花卉照片构造一个具有三类知识的训练单元。它包含原照片$\bm{x}_i$、来自同一花朵的两张视图$\bm{x}_i^{a},\bm{x}_i^{b}$，以及遮蔽后的照片$\widetilde{\bm{x}}_i$与遮蔽区域的真实像素向量$\bm{t}_i\in[0,1]^d$，其中$d>0$是该区域的像素分量数。共享表示模型$f_{\bm{w}}$连接分类头与恢复头，头部参数也计入$\bm{w}$。类别为茶梅时，分类项取$\ell_1=-\log p_{\bm{w}}(\text{茶梅}\mid\bm{x}_i)$；已确认两视图对应时，关系项取$\ell_2=\lVert f_{\bm{w}}(\bm{x}_i^{a})-f_{\bm{w}}(\bm{x}_i^{b})\rVert_2^2$，惩罚这对表示的差异。恢复头输出$\widehat{\bm{t}}_{\bm{w}}(\widetilde{\bm{x}}_i)\in\mathbb{R}^d$，恢复项取$\ell_3=\lVert\widehat{\bm{t}}_{\bm{w}}(\widetilde{\bm{x}}_i)-\bm{t}_i\rVert_2^2/d$，只在指定遮蔽区域平均像素平方误差。这里选择的距离、像素尺度和区域都是教学配置；仅拉近正配对可能使所有表示相同，不能单靠关系项保证表示有区分能力。

设该单元的茶梅概率为$0.8$，两份二维表示为$(0.2,0.4)$与$(0.3,0.2)$，两维恢复目标为$(0.6,0.2)$、预测为$(0.5,0.4)$。三项损失依次为$-\log0.8\approx0.2231$、$0.1^2+0.2^2=0.05$和$(0.1^2+0.2^2)/2=0.025$。只用这个单元，令$m_{1,r}=1$、权重向量$\bm{\alpha}=(1,0.5,2)$，各有效分母均为一，总损失约为$0.2231+0.5\times0.05+2\times0.025=0.2981$。若对应关系未经确认，就令$m_{1,2}=0$且不计算关系损失，其余两项仍可使用。固定监督、掩码与权重，在模型及损失可微、类别目标概率为正时，可以把各有效项梯度按相同权重和有效数合成，执行$\bm{w}^{(t+1)}=\bm{w}^{(t)}-\varepsilon\nabla\mathcal{L}(\bm{w}^{(t)})$，其中$\varepsilon>0$为步长；这说明三类知识怎样接入一次参数更新，不保证任意步长都降低损失。

图~\ref{fig:knowledge-fusion-heterogeneous-training}沿这三种输入追踪计算：类别标签与分类输出比较，已知视图对应与两份表示比较，被遮蔽内容与恢复输出比较。图内$m_r\alpha_r\mathcal{L}_r$简写了各分支的样本选择与贡献，不表示先计算缺失目标的损失再乘零；正文目标先筛有效单元，还包含逐项有效数归一化。三条路径共享表示参数，不共享监督语义；不同编码器的向量没有锚点映射时，也不能因维度相同就逐维相加。

\begin{figure*}[htbp]
  \centering
  \begin{tikzpicture}[
  x=0.8mm,y=1mm,
  every node/.style={font=\figurefont\small,text=BookInk},
  input/.style={book node,draw=BookBlue,fill=FigureInput!10,line width=0.9pt,minimum width=22mm,minimum height=10mm,inner sep=2pt},
  model/.style={book node,draw=BookTeal,fill=FigureModel!10,line width=1pt,minimum width=28mm,minimum height=15mm,inner sep=2pt},
  branch/.style={book node,draw=BookTeal,fill=FigureModel!7,line width=0.8pt,minimum width=27mm,minimum height=9mm,inner sep=2pt},
  knowledge/.style={book node,draw=BookAmber,fill=BookAmber!9,line width=0.8pt,minimum width=45mm,minimum height=10mm,inner sep=0pt},
  knowledge text/.style={font=\figurefont\small,anchor=west,align=left,inner sep=0pt},
  edge label/.style={font=\figurefont\scriptsize,text=BookAmber,inner sep=0pt},
  loss/.style={book node,draw=BookInk,fill=BookPaper,line width=0.8pt,minimum width=31mm,minimum height=10mm,inner sep=2pt}
]
  \node[input] (image) at (0,0) {图像与增强视图};
  \node[model] (encoder) at (37,0) {共享表示模型\\$f_{\bm{w}}$};
  \draw[book arrow,knowledge arrow,draw=BookBlue,line width=1pt] (image) -- (encoder);

  \node[branch] (pred) at (78,32) {分类预测};
  \node[branch] (compare) at (78,0) {视图表示比较};
  \node[branch] (recover) at (78,-32) {局部内容重建};
  \node[loss] (l1) at (130,32) {$m_1\alpha_1\mathcal{L}_{\mathrm{cls}}$};
  \node[loss] (l2) at (130,0) {$m_2\alpha_2\mathcal{L}_{\mathrm{rel}}$};
  \node[loss] (l3) at (130,-32) {$m_3\alpha_3\mathcal{L}_{\mathrm{rec}}$};
  \node[knowledge] (label) at (105,52) {};
  \node[knowledge] (pair) at (105,20) {};
  \node[knowledge] (patch) at (105,-52) {};
  % Reserve the left 13 mm for icons; text uses the remaining column.
  \node[knowledge text] at ([xshift=13mm]label.west) {外部任务类别标签};
  \node[knowledge text] at ([xshift=13mm]pair.west) {外部视图对应关系};
  \node[knowledge text] at ([xshift=13mm]patch.west) {外部局部恢复目标};
  \node[rectangle,draw=BookAmber,fill=white,line width=.8pt,minimum size=5mm,inner sep=0pt,font=\figurefont\bfseries\scriptsize] at ([xshift=6.5mm]label.west) {A};
  \node[circle,draw=BookAmber,fill=white,line width=.8pt,minimum size=3.5mm,inner sep=0pt] (pair-left) at ([xshift=4mm]pair.west) {};
  \node[circle,draw=BookAmber,fill=white,line width=.8pt,minimum size=3.5mm,inner sep=0pt] (pair-right) at ([xshift=9mm]pair.west) {};
  \draw[BookAmber,line width=.8pt] (pair-left) -- (pair-right);
  \node[rectangle,draw=BookAmber,fill=white,line width=.8pt,minimum size=6mm,inner sep=0pt] (patch-icon) at ([xshift=6.5mm]patch.west) {};
  \draw[BookAmber,line width=.5pt] ([xshift=-3mm]patch-icon.center) -- ([xshift=3mm]patch-icon.center);
  \draw[BookAmber,line width=.5pt] ([yshift=-3mm]patch-icon.center) -- ([yshift=3mm]patch-icon.center);
  \fill[BookAmber!45] (patch-icon.center) rectangle ([xshift=3mm,yshift=3mm]patch-icon.center);
  \node[loss,draw=BookTeal,fill=FigureModel!12,line width=1pt,minimum width=20mm] (sum) at (180,0) {总损失};

  \draw[book arrow,knowledge arrow,draw=BookTeal] (encoder.east) -- ++(3,0) |- (pred.west);
  \draw[book arrow,knowledge arrow,draw=BookTeal] (encoder.east) -- (compare.west);
  \draw[book arrow,knowledge arrow,draw=BookTeal] (encoder.east) -- ++(3,0) |- (recover.west);
  \draw[book arrow,knowledge arrow,draw=BookTeal] (pred) -- node[below,font=\figurefont\scriptsize,text=BookTeal] {预测输出} (l1);
  \draw[book arrow,knowledge arrow,draw=BookTeal] (compare) -- node[below,font=\figurefont\scriptsize,text=BookTeal] {预测输出} (l2);
  \draw[book arrow,knowledge arrow,draw=BookTeal] (recover) -- node[above,font=\figurefont\scriptsize,text=BookTeal] {预测输出} (l3);
  \draw[book arrow,knowledge arrow,draw=BookAmber] (label.south) -- ++(0,-5) -| (l1.north);
  \draw[book arrow,knowledge arrow,draw=BookAmber] (pair.south) -- ++(0,-5) -| (l2.north);
  \draw[book arrow,knowledge arrow,draw=BookAmber] (patch.north) -- ++(0,5) -| (l3.south);
  \node[edge label,anchor=south,yshift=1mm] at ($(label.south)!0.5!(l1.north)$) {监督目标};
  \node[edge label,anchor=south,yshift=1mm] at ($(pair.south)!0.5!(l2.north)$) {监督关系};
  \node[edge label,anchor=north,yshift=-1mm] at ($(patch.north)!0.5!(l3.south)$) {监督目标};
  \draw[book arrow,knowledge arrow,draw=BookTeal,line width=1pt] (l1.east) -| (sum.north);
  \draw[book arrow,knowledge arrow,draw=BookTeal,line width=1pt] (l2.east) -- (sum.west);
  \draw[book arrow,knowledge arrow,draw=BookTeal,line width=1pt] (l3.east) -| (sum.south);

  \node[anchor=north,font=\figurefont\scriptsize,text=BookMuted,align=center] at (130,-61) {掩码$m_r$在损失处选择具有相应标注的样本};
\end{tikzpicture}

  \caption{异构知识形成联合训练约束。图像经共享表示模型进入三个损失分支；任务标签、视图关系和局部恢复目标分别提供不同监督对象。掩码先选具有相应知识的单元，再计算损失，经逐项有效数归一化与权重汇总组成总目标；图内乘积是这一选择与汇总的简写，不对缺失目标求值。共享训练不要求把三类知识压缩成同一类别标签。}
  \label{fig:knowledge-fusion-heterogeneous-training}
\end{figure*}
\FloatBarrier

这些损失通过同一参数更新相互影响，改权重会改变合成梯度，若同时改采样或训练安排，还会改变资源分配。需要分别检查适用数据、尺度、采样比例和单项贡献；梯度方向长期冲突时，可以调整权重、交替批次、分阶段训练或减少共享层，但没有一种处理自动保证全部任务改善。各项仍须有独立验证指标，总损失下降不能定位哪种知识有用。

\subsection{追加来源与跨来源预算的分配}
\label{sec:knowledge-fusion-source-budget}

对$x_5$这样的规则与人工分歧，我们还可以追加判断，但最便宜的一次调用未必补充信息。同一确定模型在相同输入上再次运行可能只复制答案，独立专家却可能改变候选排序；补充规则或数据也有不同覆盖。跨来源预算要比较当前缺口下的预期信息或任务收益，而不是只按单次价格排序。

比较之前先规定当前一步视野内允许的追加行动集合$A$，包括可执行的来源调用或开发方案；当前记录与配置记为$\mathcal{D}$，行动$a$可能返回$Z_a$，$\mathcal{D}\cup Z_a$表示完成该行动后的信息状态。这里比较到该行动结束为止的收益，不自动包含后续调用或依据反馈选择下一步的价值。任务效用$U$越高越好，例如可以取既定决策损失的负期望，并固定评价对象与损失规则。结果尚未知，所以按$Z_a\mid\mathcal{D}$的预测分布平均追加后的效用，再减去当前效用，衡量期望增量。人工时间、费用、计算量与反馈时延单位不同，分别记成本$C_j(a)$和预算$B_j$，$j=1,\ldots,J$；未给换算规则时，只在每项资源允许的行动中最大化增量：
\begin{equation}
  a^\star
  \in\argmax_{a\in A}
  \left\{
    \mathbb{E}_{Z_a\mid\mathcal{D}}
    [U(\mathcal{D}\cup Z_a)]
    -U(\mathcal{D})
  \right\}
  \quad
  \text{s.t.}\quad
  C_j(a)\leq B_j,\ j=1,\ldots,J.
  \label{eq:knowledge-fusion-source-budget}
\end{equation}
目标量保留任务效用单位，约束各自保留资源单位，因此没有将小时、金额与效用直接相减。若用途要求净效用，则须先规定$\beta_j$，让$\beta_jC_j(a)$与$U$同量纲，再从期望增量减去$\sum_j\beta_jC_j(a)$。返回结果分布、收益和换算系数都需要历史审计、覆盖与依赖模型或明确业务取舍；式~\eqref{eq:knowledge-fusion-source-budget}提供比较框架，不凭空给出真实最优动作。

回到表~\ref{tab:knowledge-fusion-source-records}，在相同输入与版本上重跑确定规则$R_1$不会给$x_5$新增信息；追加独立人工判断可能帮助核对异议，新规则则需先开发，之后可能批量覆盖同类照片。预算还应计入维护、调用时延、审计和失败重试。低单次价格不等于低有效信息成本，批量复用也不能代替独立质量验证。

追加何时停止也必须沿同一口径判断。采用上述一步比较规则时，当前候选中所有预算可行动作的预期增量都不为正，或没有候选满足预算，就停止这一视野内的追加；若已给成本换算，则按同一视野内的净增量是否为正决定是否继续。这个局部规则不保证多步最优：两次调用可能相互补足，开发本身没有即时收益，也可能使后续调用获得价值。若允许联合调用、多步计划或依据反馈调整行动的策略，应把待比较方案纳入$A$，按整个方案的返回分布、效用和总成本重新比较，不能由单次增量非正排除。证据达到接收标准后进入使用，关键对象无法自动解决时转人工，预算不足而证据仍不足时保留未决并由使用端弃权。小批试用可先估计新来源的互补性，再决定扩大投入；人工内部如何安排新增与重复标注，仍由相应获取策略处理。

最后把表~\ref{tab:knowledge-fusion-source-records}的$x_2$写成可交给质量闭环的版本记录K17。该行的$R_1,R_2$同模板且报山茶，$M$报茶梅，$H$未请求，对象和物种语义已对齐；K17固定采用表~\ref{tab:knowledge-fusion-comparison}的等权相对多数，唯一最高票选类，平票或零票未决。因此山茶以$2:1$获选，却与独立参考茶梅不符。为定位这个错误，记录要保留来源原输出、共享关系、缺失原因、票数及决策配置；$2/3$仅是有效票中的山茶比例，不是后验。

K17虽然保存了共享关系，计票仍把两条规则分别当一票，记录完整并不等于规则已经利用依赖。这一行可用于诊断与选择修正，但已看过的$x_2$参考不能再次充当最终留出证据。第~\ref{chap:knowledge-quality-assessment-correction}章将接续诊断形成K18配置，并在第~\ref{sec:knowledge-quality-assessment-correction-repair-validation}节分别检查已改结果与冻结策略的独立留出质量。本章的一次错误和教学权重只揭示待核查机制，不足以决定新版本是否发布。

\section{本章小结}
\label{sec:knowledge-fusion-summary}

章首的人工与规则分歧，不能只靠增加票数解决。融合前须确认来源覆盖、错误方式与共享信息，再对齐对象和命题；类别粒度或输出形式改变时，还须说明允许转换及信息损失。来源模型可以先介绍，实际估计与推断仍须满足这些条件，票数、权重和似然才有可追踪的含义。

固定类别上，投票按票数选类，加权按得分选类，Dawid--Skene和标签模型按生成假设形成后验。无票或依赖模型增加参数与证据要求，编码不改变原动作。无参考估计须检查重叠、方向及依赖结构；完整三组联合分布可观测或可恢复时，三组秩门槛提供一类充分条件。EM后验须对应最终参数，资源停止、数值收敛、可识别与正确性则须分开。整体标签的硬约束排除组合，软约束用罚值取舍，均依赖规则适用性与求解能力。

交付取决于用途：离散行动比较错误损失，软训练保留各类概率质量，候选集合保留未区分答案，未决时使用端弃权。已确认集合事件与模型候选不能混用；异构关系和恢复目标各留相应损失，先筛有效单元再联合训练。前文案例分别表明，共享模板夸大票数，错配制造伪冲突，零有效票不增加类别证据；追加来源须针对缺口并遵守预算，一步无收益不能排除联合或后续方案的价值。保留来源、对齐、参数和使用规则的版本后，仍须独立质量评估检验可靠性。
